% =====================================================================
%  CHAPTER 10: 各类食品卫生及其管理 (Food Hygiene by Category)
% =====================================================================
\chapter{各类食品卫生及其管理}
\setcounter{chapter}{10}
\chapterintro{
\keyword{掌握}畜禽肉卫生（肉的成熟过程）；掌握油脂酸败指标（AV/POV/CGV/MDA）。\keyword{熟悉}乳和蛋类卫生；保健食品/GMF/绿色食品/有机食品。
}

% =====================================================================
%  第一节 粮豆类卫生
% =====================================================================
\section{粮豆类卫生}

\begin{keybox}
\textbf{粮豆类的主要卫生问题：}
\begin{enumerate}[leftmargin=*]
    \item \textbf{微生物污染：}霉菌及其毒素是粮豆类最重要的卫生问题（黄曲霉毒素、镰刀菌毒素、赭曲霉毒素等）。
    \item \textbf{农药残留：}田间施药和储藏期熏蒸。
    \item \textbf{其他有害化学物质的污染：}工业废水灌溉、环境重金属本底值、有毒重金属（汞、镉、铅、砷）。
    \item \textbf{仓库害虫：}温度18--21$^\circ$C、相对湿度$>$65\%时害虫繁殖最活跃。常见害虫包括谷象、米象、谷蠹、螨类等，造成数量损失和质量损害。
    \item \textbf{其他问题：}自然陈化（储存过久风味营养降低）、有毒植物种子的污染（毒麦、麦仙翁籽、毛果洋茉莉等）、无机夹杂物（砂石、泥土）、掺杂、掺假、转基因。
\end{enumerate}

\textbf{粮豆类的卫生管理措施：}
\begin{enumerate}[leftmargin=*]
    \item \textbf{控制安全水分：}粮谷安全水分为12\%--14\%，豆类为10\%--13\%。将水分控制在安全水分以下是防霉最关键的措施。
    \item \textbf{控制储藏温度：}低温储藏（$\leq$ 10$^\circ$C可有效抑制霉菌生长）。
    \item \textbf{执行真菌毒素限量标准：}GB 2761-2017规定了谷物中AFB1、DON、ZEN等限量。
    \item \textbf{加强农药残留监管：}执行GB 2763-2021规定的MRLs。
    \item \textbf{防止仓库害虫：}清洁仓库、控制温湿度、合理使用熏蒸剂。
    \item \textbf{防止无机有害物质和有毒种子的污染。}
\end{enumerate}
\end{keybox}

% =====================================================================
%  第二节 蔬菜、水果类卫生
% =====================================================================
\section{蔬菜、水果类卫生}

\begin{keybox}
\textbf{蔬菜、水果的主要卫生问题：}
\begin{enumerate}[leftmargin=*]
    \item \textbf{细菌、霉菌及其毒素污染：}人畜粪便施肥可引入肠道致病菌和寄生虫卵。腐败菌可导致蔬菜水果腐烂溃烂。
    \item \textbf{寄生虫和虫卵的污染：}未经无害化处理的人畜粪便施肥，可引入蛔虫卵、钩虫卵等。
    \item \textbf{有害化学物污染：}农药残留（有机磷、拟除虫菊酯类等）、有害重金属（工业废水灌溉）、多环芳烃类化合物（大气沉降）、硝酸盐/亚硝酸盐（过量施用氮肥）。
\end{enumerate}

\textbf{农药、硝酸盐污染的预防：}
\begin{itemize}[leftmargin=*]
    \item 人畜粪便无害化处理后使用
    \item 生食前彻底清洗
    \item 剔除烂根残叶，防止肠道致病菌和寄生虫卵的污染
    \item 选用高效低毒低残留农药，制定和执行MRLs
    \item 慎重使用激素类农药
    \item 工业废水灌溉须经无害化处理，尽量采用地下灌溉
\end{itemize}

\textbf{蔬菜、水果贮藏的卫生要求：}
\begin{itemize}[leftmargin=*]
    \item 低温冷藏（一般蔬菜0--10$^\circ$C）
    \item 控制相对湿度
    \item 气调储藏（降低O$_2$、增加CO$_2$）
    \item 避免机械损伤
\end{itemize}
\end{keybox}

% =====================================================================
%  第三节 畜禽肉类及其制品卫生
% =====================================================================
\section{畜禽肉类及其制品卫生}

\subsection{肉的成熟过程}

\begin{keybox}
\textbf{肉的成熟过程（掌握）：}

\begin{enumerate}[leftmargin=*]
    \item \textbf{僵直（Rigor）：}宰后糖原和含磷有机化合物分解为乳酸和游离磷酸，pH降至5.4--6.7。pH=5.4时达到肌凝蛋白等电点，肌凝蛋白凝固，肌纤维硬化出现僵直。
    \item \textbf{后熟（After Ripening）：}糖原继续分解产乳酸，pH回升。肌肉结缔组织变软并具弹性，此时肉松软多汁、滋味鲜美。表面因蛋白凝固形成一层干膜，可阻止微生物侵入。
    \item \textbf{自溶（Autolysis）：}常温下存放过久，组织酶在无菌条件下继续活动，分解蛋白质、脂肪。蛋白质分解。
    \item \textbf{腐败（Putrefaction）：}自溶为细菌的入侵繁殖创造条件，细菌的酶使蛋白质、含氮物质分解，pH升高。产生H$_2$S、胺类、吲哚等恶臭物质，肉色变暗绿色（硫血红蛋白形成）。
\end{enumerate}

\textbf{鉴定指标：}感官（颜色、气味、弹性和黏度）+ TVBN。
\end{keybox}

% =====================================================================
%  第二节 油脂酸败及其指标
% =====================================================================
\section{油脂酸败及其指标}

\subsection{油脂酸败（Oil Rancidity）}

\begin{keybox}
\textbf{油脂酸败（Oil Rancidity）：}食用油脂在储藏过程中由于光、热、氧、水分和金属离子等的作用，发生一系列化学变化，使其感官性状变劣和产生不良气味的现象。

\textbf{原因：}
\begin{enumerate}[leftmargin=*]
    \item \textbf{生物学原因（水解）：}脂肪酶水解甘油三酯为甘油和游离脂肪酸。
    \item \textbf{化学原因（自动氧化——最主要原因）：}不饱和脂肪酸在O$_2$作用下经自由基链式反应，产生氢过氧化物，进一步分解为醛、酮、酸等。
\end{enumerate}
\end{keybox}

\subsection{油脂酸败的理化指标}

\begin{keybox}
\textbf{油脂酸败的理化指标（掌握*）：}

\begin{enumerate}[leftmargin=*]
    \item \textbf{酸价（Acid Value, AV）：}1g油脂中游离脂肪酸所需KOH的mg数。反映油脂水解程度。

    \item \textbf{过氧化值（Peroxide Value, POV）：}100g油脂中与KI反应析出碘的g数。POV $>$ 0.25g/100g为超标。反映油脂氧化初期产生的氢过氧化物含量——是脂肪酸败最早期、最灵敏的指标。

    \item \textbf{羰基价（Carbonyl Value, CGV）：}CGV $>$ 50meq/kg为酸败。反映油脂氧化分解产生的含羰基化合物（醛、酮等）的总量——与POV形成互补。POV反映初级氧化产物，羰基价反映次级氧化（分解）产物。

    \item \textbf{丙二醛（Malondialdehyde, MDA）：}MDA $>$ 0.25mg/100g为超标。常用于反映动物油脂酸败的程度，优点是可反映甘油三酯以外的其他物质（如磷脂）的氧化破坏程度。
\end{enumerate}
\end{keybox}

% =====================================================================
%  第三节 乳及乳制品卫生
% =====================================================================
\section{乳及乳制品卫生}

\begin{defbox}
\textbf{乳及乳制品卫生：}

\begin{itemize}[leftmargin=*]
    \item \textbf{巴氏杀菌乳：}LTLT（62.8$^\circ$C、30 min）或HTST（71.7$^\circ$C、15 s），杀灭所有致病菌，需冷链配送。
    \item \textbf{灭菌乳：}UHT（135--150$^\circ$C、2 s），达到商业无菌，常温保存。
    \item \textbf{发酵乳：}以乳为原料经乳酸菌发酵制成。
\end{itemize}

\textbf{主要卫生问题：}微生物污染（沙门菌、结核杆菌、布鲁氏菌、金葡菌肠毒素）、化学性污染（农药残留、兽药残留、AF M1）、掺伪。
\end{defbox}

% =====================================================================
%  第四节 蛋类卫生
% =====================================================================
\section{蛋类卫生}

\begin{tipbox}
\textbf{蛋类卫生：}鲜蛋的变质主要是微生物通过蛋壳气孔进入蛋内引起。沙门菌是禽蛋最主要的食源性致病菌。

\textbf{鲜蛋变质过程：}
\begin{enumerate}[leftmargin=*]
    \item \textbf{贴壳蛋：}蛋白变稀，蛋黄膜破裂，蛋黄贴于蛋壳内壁。
    \item \textbf{散黄蛋：}蛋黄膜完全破裂，蛋黄与蛋白混合。
    \item \textbf{浑汤蛋/黑腐蛋：}细菌大量繁殖，蛋白质彻底分解产生H$_2$S、胺类等，有恶臭。
\end{enumerate}
\end{tipbox}

% =====================================================================
%  第五节 水产品卫生
% =====================================================================
\section{水产品（鱼虾类）卫生}

\begin{defbox}
\textbf{水产品（鱼虾类）的主要卫生问题：}
\begin{enumerate}[leftmargin=*]
    \item \textbf{腐败变质（最主要问题）：}鱼类营养丰富、水分含量高（70\%--80\%）、污染的微生物多，且组织酶的活性高。与肉类相比，更易发生腐败变质。
    \item \textbf{有害金属富集：}甲基汞在鱼贝类中经食物链逐级富集；镉可富集于贝类体内。
    \item \textbf{病原微生物：}海产食品最易受到副溶血性弧菌（嗜盐菌）污染，是夏秋季节食物中毒的重要原因。淡水鱼虾挥发性盐基总氮 $\leq$ 20 mg/100g。
    \item \textbf{寄生虫感染：}我国常见鱼类寄生虫有华支睾吸虫（肝吸虫）、肺吸虫等。生食或半生食淡水鱼虾容易感染。
    \item \textbf{天然有毒品种：}河豚含河豚毒素（TTX）；含组胺高的青皮红肉鱼类（秋刀鱼、金枪鱼、鲐鱼）可致组胺中毒。高组胺鱼类组胺限值 $\leq$ 40 mg/100g。
\end{enumerate}

\textbf{鱼类保鲜的卫生要求：}
\begin{itemize}[leftmargin=*]
    \item 海水鱼虾挥发性盐基总氮 $\leq$ 30 mg/100g
    \item 淡水鱼虾挥发性盐基总氮 $\leq$ 20 mg/100g
    \item 有效的保鲜措施：低温（冷藏/冷冻）、盐腌、防止微生物污染和减少鱼体损伤
    \item \textbf{K值：}鱼类鲜度灵敏指标。K $<$ 20\%新鲜，K $>$ 40\%为初期腐败。
\end{itemize}
\end{defbox}

% =====================================================================
%  第七节 酒类卫生
% =====================================================================
\section{酒类卫生}

\begin{keybox}
\textbf{酒类的主要卫生问题：}

\textbf{蒸馏酒与配制酒中可能存在的有害物质：}
\begin{enumerate}[leftmargin=*]
    \item \textbf{甲醇：}主要来自制酒原辅料（薯干、马铃薯、水果、糠麸等）中的果胶。具有剧烈的神经毒性，主要侵害视神经，导致视网膜受损、视神经萎缩、视力减退和双目失明。长期少量摄入可导致慢性中毒，特征性临床表现为视野缩小、不可校正的视力减退。
    \item \textbf{醛类：}如乙醛、糠醛，为酒醅发酵过程中生成的中间产物。乙醛毒性大于甲醛。
    \item \textbf{杂醇油：}由蛋白质、氨基酸和糖类分解产生，包括丙醇、异丁醇、异戊醇等。毒性及麻醉力比乙醇强，使中枢神经系统充血，引起剧烈头痛。
    \item \textbf{氢氰酸：}来自木薯或其他含氰苷原料。中毒使机体陷于窒息状态，麻痹呼吸中枢及血管运动中枢。
    \item \textbf{铅：}来自蒸馏器、冷凝导管和储酒容器中的铅溶出。
    \item \textbf{锰：}高锰酸钾-活性炭脱臭除杂处理时引入。
\end{enumerate}

\textbf{发酵酒中的卫生问题：}展青霉素（水果原料霉变）、二氧化硫残留、亚硝胺、微生物污染（细菌、酵母菌）、添加剂使用问题。
\end{keybox}

% =====================================================================
%  第八节 冷饮食品卫生
% =====================================================================
\section{冷饮食品卫生}

\begin{keybox}
\textbf{冷饮食品的分类：}冰淇淋、雪糕、冰棍、食用冰、碳酸饮料、果汁饮料、乳酸菌饮料等。

\textbf{主要卫生问题：}
\begin{enumerate}[leftmargin=*]
    \item \textbf{微生物污染：}原料含菌量高、加工过程污染、杀菌不彻底、冷链不完善。致病菌包括沙门菌、金黄色葡萄球菌肠毒素等。
    \item \textbf{食品添加剂：}色素、香精、甜味剂、酸度调节剂等需符合GB 2760规定。
    \item \textbf{重金属污染：}铅、铜、砷等来自加工设备和原料。
\end{enumerate}
\end{keybox}

% =====================================================================
%  第九节 罐头食品卫生
% =====================================================================
\section{罐头食品卫生}

\begin{keybox}
\textbf{罐头食品的分类：}按内容物分为肉类罐头、禽类罐头、水产品罐头、水果蔬菜罐头等；按包装分为金属罐、玻璃罐、复合袋（软罐头）等。

\textbf{主要卫生问题：}
\begin{enumerate}[leftmargin=*]
    \item \textbf{罐头容器的卫生要求：}无渗漏、无锈蚀、内壁涂料完整无脱落。
    \item \textbf{原料微生物污染与灭菌：}必须达到商业无菌。D值、F值（通常12D）概念应用于灭菌条件设定。
    \item \textbf{重金属污染：}酸性内容物可腐蚀金属罐、产生氢气、溶出锡和铅。
\end{enumerate}

\textbf{罐头食品的卫生学鉴定及处理：}
\begin{enumerate}[leftmargin=*]
    \item \textbf{感官检查：}容器密封完好、无泄漏、无胖听；内容物具有该品种应有的色泽、气味、滋味、形态。
    \item \textbf{胖听（Swelling）：}由于罐内微生物活动或化学作用产生气体，形成正压，使一端或两端外凸。
    \begin{itemize}
        \item \textbf{化学性胖听：}金属罐受酸性内容物腐蚀产生氢气，叩击呈鼓音，穿洞有气体逸出但无腐败气味，一般不宜食用。
        \item \textbf{生物性胖听：}微生物活动产气所致，叩击呈浊音，穿洞有腐败臭味，\textbf{绝对不可食用！}
        \item \textbf{物理性胖听：}低温装罐、运输海拔变化等物理原因导致，内容物未变质，可食用。
    \end{itemize}
    \item \textbf{平酸腐败（Flat-sour Spoilage）：}由分解碳水化合物产酸不产气的平酸菌引起，内容物酸度增加而外观完全正常。
\end{enumerate}
\end{keybox}

% =====================================================================
%  第十节 调味品卫生
% =====================================================================
\section{调味品卫生}

\begin{keybox}
\textbf{调味品的定义和营养学意义：}以粮食、豆类、水果等为原料，经发酵或调配制成的具有调味功能的产品。主要提供咸、甜、酸、鲜、辣等基本味道，酱油和发酵酱提供少量蛋白质和B族维生素，食盐提供钠，食醋提供有机酸。

\textbf{酱油的主要卫生问题：}
\begin{itemize}[leftmargin=*]
    \item 原料霉变导致黄曲霉毒素污染
    \item 氯丙醇（3-MCPD）——酸水解植物蛋白液（HVP）配制酱油中的有害物
    \item 微生物污染（耐盐酵母菌、细菌等产生霉变和"生白"）
    \item 防腐剂使用
    \item 掺伪和假冒
\end{itemize}

\textbf{食醋的卫生问题：}
\begin{itemize}[leftmargin=*]
    \item 原料卫生
    \item 醋酸发酵过程的微生物管理
    \item 不应使用冰醋酸勾兑
    \item 昆虫（醋鳗、醋蝇）污染
\end{itemize}

\textbf{食盐的卫生问题：}
\begin{itemize}[leftmargin=*]
    \item 碘缺乏病——我国实行食盐加碘（GB 26878）
    \item 杂质控制（不溶性杂质、重金属）
    \item 亚铁氰化钾等抗结剂的使用
\end{itemize}
\end{keybox}

% =====================================================================
%  第十一节 特殊食品
% =====================================================================
\section{特殊食品}

\subsection{保健食品（Health Food）}

\begin{defbox}
\textbf{保健食品（Health Food）：}声称具有特定保健功能或以补充维生素、矿物质为目的的食品。即适用于特定人群食用，具有调节机体功能，不以治疗疾病为目的，且对人体不产生任何急性、亚急性或者慢性危害的食品。
\end{defbox}

\subsection{转基因食品（GMF）}

\begin{tipbox}
\textbf{转基因食品（Genetically Modified Food, GMF）：}利用基因工程技术改变基因组构成的生物所生产的食品。

\textbf{GMF分类：}分为I/II/III/IV类。安全性评价遵循实质等同性原则和个案评价原则。
\end{tipbox}

\subsection{无公害食品（Non-environmental Pollution Food）}

\begin{defbox}
\textbf{无公害食品（熟悉）：}产地环境、生产过程和产品质量符合国家有关标准和规范的要求，经认证合格获得认证证书并允许使用无公害农产品标志的、未经加工或者初加工的食用农产品。无公害食品是食品安全的基本门槛。
\end{defbox}

\subsection{绿色食品（Green Food）}

\begin{defbox}
\textbf{绿色食品（Green Food，熟悉）：}遵循可持续发展原则，按照绿色食品标准生产，经过专门机构认定，许可使用绿色食品标志的无污染、安全、优质、营养类食品。具备安全和营养的双重质量保证。
\begin{itemize}[leftmargin=*]
    \item \textbf{A级：}限量使用化学合成生产资料。
    \item \textbf{AA级：}完全按照有机生产方式生产（等同于有机食品）。
\end{itemize}
\end{defbox}

\subsection{有机食品（Organic Food）}

\begin{defbox}
\textbf{有机食品（Organic Food，熟悉）：}来自于有机农业生产体系，根据有机农业生产的规范生产加工，并经独立的认证机构认证的农产品及其加工产品。国际通行的最高级别认证，禁止使用任何化学合成的农药、化肥、激素、抗生素、转基因技术。
\end{defbox}

% =====================================================================
%  复习题
% =====================================================================
\reviewcontent{
\section{复习题}

\subsection*{一、名词解释}

\begin{enumerate}[leftmargin=*, itemsep=6pt]
    \item \textbf{僵直（Rigor）、后熟（After Ripening）、自溶（Autolysis）、腐败（Putrefaction）}
    \item \textbf{油脂酸败、酸价（AV）、过氧化值（POV）、羰基价（CGV）、丙二醛（MDA）}
    \item \textbf{巴氏杀菌乳、灭菌乳、K值}
    \item \textbf{保健食品、转基因食品（GMF）、绿色食品、有机食品}
\end{enumerate}

\subsection*{二、简答题}

\begin{enumerate}[leftmargin=*, itemsep=8pt]
    \item \textbf{【重点】}简述畜禽肉成熟的四个阶段及各阶段的特征。
    \item \textbf{【重点】}简述油脂酸败的定义、原因及四大理化指标（AV、POV、CGV、MDA）的卫生学意义。
    \item 简述乳及乳制品的种类和主要卫生问题。
    \item 简述鲜蛋变质的过程。
    \item 简述水产品的主要卫生问题及鲜度指标。
    \item 简述保健食品、转基因食品、绿色食品和有机食品的定义及区别。
\end{enumerate}

\subsection*{参考答案要点}

\subsubsection*{名词解释（要点）}

\begin{enumerate}[leftmargin=*, itemsep=4pt]
    \item \textbf{僵直：}宰后pH降至5.4--6.7，肌凝蛋白凝固，肌纤维硬化。\textbf{后熟：}pH回升，肉松软多汁滋味鲜美。\textbf{自溶：}组织酶分解蛋白质、脂肪。\textbf{腐败：}细菌作用，pH升高，产生H$_2$S、胺类、吲哚等恶臭物质。
    \item \textbf{油脂酸败：}油脂在光、热、氧、水分和金属离子等作用下发生感官变劣和产生不良气味的化学变化。\textbf{AV：}1g油脂中游离脂肪酸所需KOH的mg数。\textbf{POV：}100g油脂中与KI反应析出碘的g数，POV $>$ 0.25g/100g为超标，氧化早期最灵敏指标。\textbf{CGV：}CGV $>$ 50meq/kg为酸败，反映次级氧化产物。\textbf{MDA：}MDA $>$ 0.25mg/100g为超标。
    \item \textbf{巴氏杀菌乳：}LTLT或HTST消毒，杀灭所有致病菌。\textbf{灭菌乳：}UHT灭菌，商业无菌，常温保存。\textbf{K值：}鱼类鲜度指标，K $<$ 20\%新鲜，K $>$ 40\%腐败。
    \item \textbf{保健食品：}声称特定保健功能或以补充维生素、矿物质为目的的食品。\textbf{GMF：}利用基因工程技术改变基因组构成的生物所生产的食品，分为I/II/III/IV类。\textbf{绿色食品：}AA级和A级。\textbf{有机食品：}禁止化学合成农药、化肥、激素、抗生素、转基因技术的最高级别认证。
\end{enumerate}

\subsubsection*{简答题（要点）}

\begin{enumerate}[leftmargin=*, itemsep=6pt]
    \item \textbf{四阶段：}①僵直（pH 5.4--6.7，肌凝蛋白凝固，肌纤维硬化）；②后熟（pH回升，肉变松软多汁）；③自溶（蛋白质分解）；④腐败（pH升高，产生恶臭物质）。
    \item \textbf{原因：}生物学水解（脂肪酶）和化学自动氧化（自由基链式反应）。\textbf{AV：}反映水解程度。\textbf{POV：}氧化初期最灵敏指标，POV $>$ 0.25g/100g为超标。\textbf{CGV：}反映次级氧化产物，CGV $>$ 50meq/kg为酸败。\textbf{MDA：}反映动物油脂酸败程度，MDA $>$ 0.25mg/100g为超标。
    \item \textbf{种类：}巴氏杀菌乳、灭菌乳、发酵乳。\textbf{卫生问题：}微生物污染（沙门菌、结核杆菌、布鲁氏菌）、化学性污染（农药/兽药/AF M1）、掺伪。
    \item \textbf{变质过程：}贴壳蛋 $\rightarrow$ 散黄蛋 $\rightarrow$ 浑汤蛋/黑腐蛋。主要机制是微生物通过蛋壳气孔进入蛋内引起。
    \item \textbf{问题：}腐败变质（含水高、组织脆）、有害金属富集、病原微生物（副溶血性弧菌）、寄生虫（肝吸虫）、天然有毒（河豚）。\textbf{K值：}鱼类鲜度指标；\textbf{组胺中毒：}与青皮红肉鱼有关。
    \item \textbf{保健食品：}声称保健功能，不以治疗为目的。\textbf{GMF：}基因工程技术改造，分I/II/III/IV类。\textbf{绿色食品：}AA级（等同于有机）和A级（限量化学合成）。\textbf{有机食品：}最高认证，全禁化学合成物+转基因。
\end{enumerate}

\newpage
}
